\documentclass[12pt]{article}
\usepackage[utf8]{inputenc}
\usepackage[T1]{fontenc}
\usepackage[a4paper,margin=2.5cm]{geometry}
\usepackage{enumitem}
\setlist[itemize]{noitemsep,topsep=4pt}
\title{Sprint 3 - S3\_18\\APIs REST e Testes de Integração}
\author{Matheus Moreli Durão Silva\\Matrícula 12621CBS217\\Cibersegurança - Bacharelado - Noturno}
\date{2026}
\begin{document}
\maketitle

\section*{Entregável}
Foi produzida uma collection Postman para a API de ativos, acompanhada de guia de uso e evidência de execução. A collection possui operações GET, POST, PUT e DELETE sobre o recurso \texttt{/ativos}.

\section*{O que foi estudado e aplicado}
\begin{itemize}
\item Métodos HTTP: GET para consulta, POST para criação, PUT para atualização e DELETE para remoção.
\item Códigos de status: 200, 201 e 204 confirmam o resultado esperado de cada operação.
\item JSON: o corpo das requisições envia nome, tipo e responsável do ativo.
\item Testes de integração: cada requisição possui uma assertion de status.
\end{itemize}

\section*{Execução comprovada}
Em 13/10/2026, a collection foi executada pelo Newman 11.16.0, executor de collections do Postman. Foram executadas quatro requisições, quatro assertions aprovadas e zero falhas. O POST guarda o identificador retornado em \texttt{ativoId}; o PUT e o DELETE reutilizam esse identificador para validar o CRUD encadeado.

\section*{Como conferir}
Inicie a API da S3\_19 e execute \texttt{postman\_collection.json}. O guia \texttt{como\_usar\_collection.md} descreve o fluxo; \texttt{resultado\_newman.md} registra os códigos obtidos e o comando de execução.

\section*{Conclusão}
O entregável documenta e testa a API REST do Inventário de Ativos com uma collection reproduzível e evidência real de execução.
\end{document}
